\section{考点总结}
\begin{longtable}{P{3cm}P{5cm}P{6.1cm}}
\toprule 考点 & 实现方法 & 对应真题与常见错误\\\midrule\endhead
拆数位与基本循环 & 取余、整除、条件分支 & 数位之和、报数；检查0、跳过计数与轮次。\\
统计与值域数组 & 次数桶、边读边更新 & 门禁系统、灰度直方图；容量由值域决定。\\
二维数组与坐标映射 & 按目标输出顺序访问原数组 & 图像旋转；逆时针方向和行列交换。\\
局部关系判断 & 比较相邻项、首尾单独处理 & 折点计数、卖菜；端点不能访问越界。\\
对象与多条件判断 & 结构体、逐项代入 & 线性分类器；先确定符号与类别的对应。\\
枚举候选与最小编号 & 按编号升序，遇到首个合法者停止 & 仓库规划；严格比较，不能只比较部分维度。\\
小模拟 & 定义状态时刻和更新顺序 & 网格同步更新、日期计算；不要混用新旧状态。\\
浮点数 & 均值、方差、数学函数 & 归一化处理；整数除法和输出精度。\\
\bottomrule\end{longtable}
\textbf{学习后的自测：}能否用一句话说明变量的含义？数组最后一个合法下标是什么？最坏情况总共比较多少次？哪个中间值最大？相等、0、负值或最小规模时程序仍按定义工作吗？
